\documentclass[12pt,a4paper]{report}
\usepackage[]{amsmath,amsthm,amssymb,amscd}
\usepackage[latin1]{inputenc}
\usepackage{color}
%\usepackage{fullpage}

\topmargin 0pt
\advance \topmargin by -\headheight
\advance \topmargin by -\headsep
     
\textheight 246mm
     
\oddsidemargin 0pt
\evensidemargin \oddsidemargin
\marginparwidth 0.5in
     
\textwidth 159mm


%               Theorem Environments
\theoremstyle{definition}
\newtheorem{ex}{}
\newcommand{\pd}[1]{\frac{\partial}{\partial #1}} %\pd{x}
%               Subquestions numbered with letters
\renewcommand{\theenumi}{\textbf{\alph{enumi}}}

%               Maths definitions

\newcommand{\NN}{\ensuremath{\mathbb N}}
\newcommand{\ZZ}{\ensuremath{\mathbb Z}}
\newcommand{\CC}{\ensuremath{\mathbb C}}
\newcommand{\TT}{\ensuremath{\mathbb T}}
\newcommand{\RR}{\ensuremath{\mathbb R}}
\newcommand{\QQ}{\ensuremath{\mathbb Q}}
\newcommand{\g}{\ensuremath{\frak{g}}}
\newcommand{\h}{\ensuremath{\frak{h}}}
\newcommand{\ft}{\ensuremath{\frak{t}}}
\newcommand{\cB}{\mathcal{B}}
\newcommand{\cN}{\mathcal{N}}
\newcommand{\cL}{\mathcal{L}}
\newcommand{\cD}{\mathcal{D}}
%\newcommand{\pd}[1]{\frac{\partial}{\partial #1}} %\pd{x}
%\newcommand{\pd}[1]{\frac{\partial}{\partial #1}} %\pd{x}


\pagestyle{empty} %%%
\begin{document}
\noindent
{\sffamily\large  
Marco Zambon\hfill     30/11/2012\\\'Algebra I, curso 2012-13}
 

%\noindent
%\hrulefill{}\\
\begin{center}\bfseries\large\textsf{Ejercicios resueltos en clase 2    \vspace{-2mm}}
\end{center}

\textcolor{red}{\bfseries\Huge\textsf{Incluye solamente una versi\'on preliminar del problema 3 }}
 
%\hrulefill{}\\

\noindent
\hrulefill{}\\

\noindent
{\sffamily\large Apellidos:\\
Nombre:\\
DNI: 
}


 
\begin{center}{\sffamily (deja esta tabla vac\'ia)}
\end{center}
\vspace*{-5mm}
%\begin{table}[htdp]
%\caption{}
%\begin{center}
%\begin{tabular}{|c|c|c|c|c|c||c|}
%\hline
%Ej. 1  &Ej. 2&Ej. 3&Ej. 4&Ej. 5&Ej. 6& Total puntos (de 100)\\
%\hline
%& & & & & &  \\
%& & & & & &  \\
%\hline
%\end{tabular}
%\end{center}
%\label{default}
%\end{table}
 

\begin{table}[htdp]
%\caption{}
\begin{center}
\begin{tabular}{|c|p{2cm} |}
\hline
&\\[-0.2cm]
Ej. 1  &\\[0.2cm]
\hline
&\\[-0.2cm]
Ej. 2  &\\[0.2cm]
\hline
&\\[-0.2cm]
Ej. 3  &\\[0.2cm]
\hline
%Ej. 4  &\\[0.2cm]
%\hline
%&\\[-0.2cm]
%Ej. 5  &\\[0.2cm]
%\hline
%&\\[-0.2cm]
%Ej. 6  &\\[0.2cm]
%\hline
\hline
&\\[-0.2cm]
Total (de 15) & \\[0.2cm]
\hline
\end{tabular}
\end{center}
\label{default}
\end{table}
\vspace*{-5mm}


%%
\noindent
%\hrulefill{}\\
\textsf{{\bf \sffamily  Instrucciones:} 
%Utiliza una hoja distinta para 
%cada problema, y en cada hoja escribe  tu nombre y el n\'umero del problema.\\
A continuaci\'on encontrar\'as los enunciados de los ejercicios, con espacio debajo de ellos para escribir tu soluci\'on.\\
Justifica todas tus respuestas.
Puedes utilizar, cit\'andolos de manera apropriada, los resultados que hemos obtenido en clase.\\
Entrega unicamente estas hojas.
}
% o en los ejercicios para casa.  
 
 \noindent
\hrulefill{}\\


\newpage 
\hrulefill{}\\ 
\begin{ex}[5 puntos]   
a)  Considera la aplicaci\'on lineal de $\RR^2$ a $\RR^2$ dada por reflexi\'on con respecto a la recta $x_1=x_2$. Encuentra la matriz correspondiente.

b) Encuentra la matriz correspondiente a la aplicaci\'on lineal
$T \colon  \RR^2 \to \RR^3$,
$$T\left(\begin{array}{c}x\\y \end{array}\right) = \left(\begin{array}{c}x + 2y\\ 2x - 5y\\ 7y\end{array}\right).$$ \end{ex}  
%3.3 trail ch 1

\newpage  
\hrulefill{}\\ 
\begin{ex}[5 puntos]  Averiguar si son subespacios vectoriales de $\mathcal{M}_{2\times 2}(\mathbb{R})$ (las matrices reales 2 por 2) los siguientes subconjuntos: 
 
a) Las matrices de n\'umeros reales de orden $2\times 2$ de rango $1$. 

b) Las matrices de n\'umeros reales de orden $2\times 2$ que conmutan con la matriz 
B, siendo B una matriz fija $2\times 2$. 
\end{ex} 
%7 hoja 5

%%
\newpage 
\hrulefill{}\\ 
\begin{ex}[5 puntos]   
%Siendo  
%$$S_1 \equiv\left\{\begin{array}{c@{\hspace{3pt}}c@{\hspace{3pt}}c@{\hspace{3pt}}c@{\hspace{3pt}}c@{\hspace{3pt}}c} 
%x_1 && -x_3 &&= 0\\ 
%& x_2 & &-x_4&  = 0\end{array}\right.\qquad 
%y \qquad 
%S_2 \equiv\left\{\begin{array}{c@{\hspace{3pt}}c@{\hspace{3pt}}c@{\hspace{3pt}}c@{\hspace{3pt}}c@{\hspace{3pt}}c} 
%x_1 &+x_2 &+x_3 &+x_4 &= 0\\ 
%x_1 &+x_2&&& = 0\end{array}\right.$$
%hallar una base y las ecuaciones cartesianas de $S_1 + S_2$ y $S_1\cap S_2$.\end{ex}
Hallar las ecuaciones cartesianas del siguiente subespacio de $\mathbb{R}^4$: 
$$\begin{array}{c@{\hspace{2pt}}l}
%S_1 =& \operatorname{Span}\{(1, 0, 1, 0),(2, 1, 0, -1)(1, -1, 3, 1)\},\\[+2pt]
S := & \operatorname{Span}\{(3, 1, 0, -1),(1, 1, -1, -1), (7, 1, 2, -1)\}.\\[+2pt]
%S_3 = & \operatorname{Span}\{(0, 2, 5, 0)\}.
\end{array}$$
\end{ex}

{\bf M\'etodo 1}
Estamos buscandos $a,b,c,d\in \RR$ tales que $$S\subset \{(x,y,z,w):ax+by+cz+dw=0\}.$$ En terminos del sistema de generadores de $S$ dado arriba,  eso es equivalente  a pedir que $a,b,c,d$ satisfagan
 \begin{align*}
a\cdot 3+b\cdot 1+c\cdot 0+d\cdot(-1)=0\\
a\cdot 1+b\cdot 1+c\cdot (-1)+d\cdot(-1)=0\\
a\cdot 7+b\cdot 1+c\cdot 2+d\cdot(-1)=0.
\end{align*}
Resolvemos este sistema con el m\'etodo de Gauss: la forma escalonada de la matriz 
$$\left(\begin{array}{cccc}3 & 1 & 0 & -1 \\1 & 1 & -1 & -1 \\7 & 1 & 2 & -1\end{array}\right)$$ es
$$\left(\begin{array}{cccc}1 & 1 & -1 & -1 \\0 & -2 & 3 & 2 \\0 & 0 & 0 & 0\end{array}\right),$$
por lo tanto las soluciones del sistema arriba son exactamente 
$$\left\{c
\left(\begin{array}{c}-\frac{1}{2} \\\frac{3}{2} \\1 \\0\end{array}\right)
+
d
\left(\begin{array}{c}0 \\1 \\0 \\1\end{array}\right):c,d\in \RR\right\}.$$
Por lo tanto un par de ecuaciones cartesianas para $S$ es dado por 
$ax+by+cz+dw=0$ con $(a,b,c,d)$ igual a los dos vectores de arriba,
es decir,
$-\frac{1}{2}x+\frac{3}{2}y+z=0$ y $y+w=0$.

\bigskip

{\bf M\'etodo 2}


Sea $(x,y,z,w)$ un elemento de $S$. En otras palavras $(x,y,z,w)$ se puede escribir 
como combinaci\'on lineal de los 3 generadores dados de $S$, o equivalentemente, el sistema
$$c_1
\left(\begin{array}{c}3 \\1 \\0 \\-1\end{array}\right)
+
c_3
\left(\begin{array}{c}1 \\1 \\-1 \\-1\end{array}\right)+
c_3
\left(\begin{array}{c}7 \\1 \\2 \\-1\end{array}\right)
=
\left(\begin{array}{c}x \\y \\z \\w\end{array}\right)
$$
(con incognitas $c_1,c_2,c_3$)
es compatible. Equivalentemente,
la forma escalonada de la matriz
$$
\left(\begin{array}{ccc|c}3 & 1 & 7 & x\\1 & 1 & 1 & y \\0 & -1 & 2 & z \\-1 & -1 & -1 & w\end{array}\right)$$
no tiene ning\'una fila de la forma $(0\;0\;0|\;\beta)$ con $\beta\neq 0$.
La forma escalonada de la matriz arriba es
$$
\left(\begin{array}{ccc|c}1 & 1 & 1 & y \\
0 & -2 & 4 & x-3y\\
0 & 0 & 0 & -\frac{1}{2}x+\frac{3}{2}y+z \\0 & 0 & 0 & y+w\end{array}\right).$$
Concluimos que $(x,y,z,w)$ un elemento de $S$ si y solo si
$-\frac{1}{2}x+\frac{3}{2}y+z=0$, $y+w=0$. Por lo tanto estas son las ecuaciones cartesianas de $S$.

% 17 de hoja 6

  
 


 \end{document}
 %%%%%%%
 \begin{ex}

\end{ex}

 \begin{ex}
\begin{enumerate}
\item 
\end{enumerate}
\end{ex}
 
 
 
 